% !TeX root = ../traffic_rules.tex

\section{Predicates and Functions}
\label{sec:predicates}

The predicates and functions necessary to describe the presented traffic rules are categorized in position, velocity, braking, and general elements.
We use first-order logic to define the predicates and functions (where the atomic propositions of the FLTL formulas correspond to the defined predicates).

\subsection{Position Elements}
The occupied lanelets of a state $\kState$ are defined by the intersection of its occupancy with the corresponding lanelets:
\begin{align*}
\occupiedLaneletsFunction(\kState) = \big\{l \in \mathbb{L}   \, \big| \, l \cap \mathcal{O}(\kState) \neq \emptyset \big\}.
\end{align*}
The lanes which a vehicle occupies are defined as:
\begin{align*}
\begin{split}
\lanes(\kState) = \{\lane(l)  \, \big| \, l \in(\occupiedLaneletsFunction(\kState))\}. 
\end{split}
\end{align*}
This information is used to reason whether the $\text{k}^\mathrm{th}$ vehicle is in any lane the $\text{p}^\mathrm{th}$ vehicle is in:
\begin{align*}
\begin{split}
\inSameLaneAsEgoVehiclePredicate(\kState, \pState) \iff \lanes(\kState) \cap \lanes(\pState) \neq \emptyset. 
%& \big\{l \, \big| \, \forall l' \in \occupiedLaneletsFunction(\egoState): l \in \lane(l') \big\} \, \cap \\
%&  \big\{l \, \big| \,\forall l' \in \occupiedLaneletsFunction(\otherState):  l \in \lane(l') \big\}
% \neq \emptyset.
\end{split}
\end{align*}
The $\text{k}^\mathrm{th}$ vehicle is in front of the $\text{p}^\mathrm{th}$ vehicle if its longitudinal position is larger than the one of the $\text{p}^\mathrm{th}$ vehicle:
\begin{align*}
\inFrontOfEgoVehiclePredicate(\pState, \kState) \iff \frontPositionFunction(\pState) \leq \rearPositionFunction(\kState).
\end{align*}
The information whether a vehicle is on a the main carriage way, is given by
\begin{align*}
& \onMainCarriageWayPredicate(\kState) \iff \\
& \qquad main\_carriage\_way \in\\ 
& \qquad \quad \{\laneltAttribute(l)  \, \big| \, l \in \occupiedLaneletsFunction(\kState)\}.
\end{align*}
%where
%\begin{align*}
%& \onLaneletType(\kState) = \\
%& \qquad \big\{a \in \mathbb{A}   \, \big| \, \forall l \in \occupiedLaneletsFunction(\kState): a = attr(l) \big\}.
%\end{align*}
Other lanelet attributes can be evaluated similarly.
The $\text{k}^\mathrm{th}$ vehicle is left of the $\text{p}^\mathrm{th}$ vehicle if
\begin{align*}
\begin{split}
& \vehicleIsLeftPredicate(\kState, \pState) \iff \\
&\qquad \big(\rearPositionFunction(\pState) \leq  \rearPositionFunction(\kState) \leq  \frontPositionFunction(\pState) \\
& \qquad \lor (\frontPositionFunction(\pState) < \frontPositionFunction(\kState) \land   \rearPositionFunction(\kState) < \rearPositionFunction(\pState)) \\
& \qquad \lor \rearPositionFunction(\pState) \leq  \frontPositionFunction(\kState) \leq  \frontPositionFunction(\pState) \big) \land d_\mathrm{k} > d_\mathrm{p}, \\
\end{split}
\end{align*} 
evaluates to true. To determine whether a vehicle is right of a broad lane marking, the following predicate is introduced:
%\begin{align*}
%\begin{split}
%&  \rightOfBroadLaneMarkingPredicate(\kState) \iff \\ 
%& \quad \Big(\big\{l \in \occupiedLaneletsFunction(\kStatek) \, \big| \, \laneMarkingLeftFuction(l) = broad\_solid \\
%& \qquad \lor \laneMarkingLeftFuction(l) = broad\_dashed \big\} \neq \emptyset \, \land\\
%& \qquad  \big\{l \in \occupiedLaneletsFunction(\kState) \, \big| \, \laneMarkingRightFuction(l) = broad\_solid \\
%&\qquad  \lor \laneMarkingRightFuction(l) = broad\_dashed \big\} = \emptyset \Big) \, \lor\\
%&\quad   \Big(\exists l \in \allLaneletsLeftOfStateFunction(\kState): \\ 
%& \qquad \laneMarkingLeftFuction(l) = broad\_solid \lor \laneMarkingLeftFuction(l) = broad\_dashed\Big)
%\end{split}
%\end{align*}
\begin{align*}
\begin{split}
&  \rightOfBroadLaneMarkingPredicate(\kState) \iff \\ 
& ~~ \big\{l \in \occupiedLaneletsFunction(\kState) \, \big| \, \laneMarkingRightFuction(l) \in \{broad\_solid,  \\
& \qquad broad\_dashed\}\} = \emptyset \, \land\\
& \quad \big\{l \in \allLaneletsLeftOfStateFunction(\kState) \, \big| \, \\ 
& \qquad \laneMarkingRightFuction(l) \in \{broad\_solid, broad\_dashed\}\} \neq \emptyset, \\
\end{split}
\end{align*}
where
\begin{align*}
\begin{split}
& \allLaneletsLeftOfStateFunction(\kState) = \\ 
& \bigcup \{\allLaneletsLeftOfLaneletFunction(l) \, \big| \, l \in \occupiedLaneletsFunction(\kState)\}, %\setminus  \occupiedLaneletsFunction(\kState), 
\end{split}
\end{align*}
returns the set containing all lanelets left of the vehicle and the recursively defined function
\begin{align*}
\begin{split}
& \allLaneletsLeftOfLaneletFunction(l) = \\ 
& \qquad \{\laneletLeft(l)\} \cup \allLaneletsLeftOfLaneletFunction(\laneletLeft(l)), \\
%& \qquad \big\{l'  \, \big| \, l' \in (\laneletLeft(l) \cup \allLaneletsLeftOfLaneletFunction(\laneletLeft(l))) \big\} .
\end{split}
\end{align*}
returns all lanelets left of a given lanelet, where $\allLaneletsLeftOfLaneletFunction(\bot) = \emptyset$.
%The first clause of $\rightOfBroadLaneMarkingPredicate$ checks if the vehicle's occupied lanelets contain a broad lane marking where the vehicle is right of. The second clause checks for lanelets on the left side of the  vehicle's occupied lanelets. 
Similarly, the predicates $\leftOfBroadLaneMarkingPredicate$, $\allLaneletsRightOfStateFunction$, and $\allLaneletsRightOfLaneletFunction$ can be defined.
The width of a road given a lanelet and longitudinal position $s_\mathrm{k}$ is recursively defined as:
\begin{align*}
\begin{split}
\roadWidth(l, s_\mathrm{k}) = \sum_{l' \in \adjacenLanelets(l)} \laneletWidth(l', s_\mathrm{k}) 
\end{split}
\end{align*} 
where $\roadWidth(\bot, s_\mathrm{k}) = 0$ and 
\begin{align*}
\begin{split}
\adjacenLanelets(l) = \big(\{l\} &\cup \adjacenLanelets(\laneletRight(l)) \\ 
& \cup \adjacenLanelets(\laneletLeft(l))\big) \setminus \{\bot\}.
\end{split}
\end{align*} 
If the width of a road is larger than a certain threshold $\minLaneWidth$, the road is declared as broad enough:
\begin{align*}
\begin{split}
& \interstateBroadEnoughForEmergencyLanePredicate(\kState) \iff  \\
& \quad \forall l \in \occupiedLaneletsFunction(\kState): \roadWidth(l, s_\mathrm{k}) > \minLaneWidth. \\  
\end{split}
\end{align*}
A vehicle is in the leftmost lane if any of its occupied lanelets has no left adjacent lanelet:
\begin{align*}
\begin{split}
& \inLeftMostLanePredicate(\kState) \iff \\ 
& \qquad \exists l \in \occupiedLaneletsFunction(\kState): \laneletLeft(l) = \bot.
\end{split}
\end{align*} 
The proposition $\inRightMostLanePredicate$ for the rightmost lane is defined analogously.
To evaluate if a vehicle drives as leftmost or as rightmost as possible, the lateral distance from the lane center and the lane width are used:
\begin{align*}
\begin{split}
& \drivesLeftMostPredicate(\kState, s_\mathrm{k}) \iff \forall l \in \occupiedLaneletsFunction(\kState): \\ 
& \qquad \bigg(\frac{1}{2} \laneletWidth(l, s_\mathrm{k}) + \leftPositionFunction(\kState)\bigg) \leq \nearBoundDrivingThreshold,
\end{split}
\end{align*}
where $\drivesRightMostPredicate$ is analogously defined.


\subsection{Velocity Elements}
The speed limit regulated by traffic signs is the minimum speed limit of all lanelets a vehicle occupies:
\begin{align*}
\begin{split}
\speedLimitLane(\kState) = \min\big(\speedLimitFunction(\mathtt{\occupiedLaneletsFunction}(\kState)) \big).
%\speedLimitLane(\kState) = \min \big( \big\{&v_\mathrm{sl}   \, \big| \, \forall l \in \mathtt{\occupiedLaneletsFunction}(\kState): \\
%&v_\mathrm{sl} = \speedLimitFunction(l) \big\} \big)
\end{split}
\end{align*}
Similar, the minimum speed limit is defined as:
\begin{align*}
\begin{split}
\requiredSpeedLane(\kState) = \max\big(\requiredSpeedFunction(\mathtt{\occupiedLaneletsFunction}(\kState))\big).
%\requiredSpeedLane(\kState) = \max \big( \big\{&v_\mathrm{sl}   \, \big| \, \forall l \in \mathtt{\occupiedLaneletsFunction}(\kState): \\
%&v_\mathrm{sl} = \requiredSpeedFunction(l) \big\} \big)
\end{split}
\end{align*}
The ego vehicle has to be able to stop within the field of view. Therefore, the maximum allowed velocity based on the field of view \speedLimitFOV is chosen such that the ego vehicle can come safely to a standstill even when a standing vehicle enters the field of view $fov$. For this, we search for the maximum value of \speedLimitFOV which fulfills the following condition: 
\begin{align*}
\begin{split}
& \frontPositionFunction(\xi(t_{s}(u_\mathrm{brake}(\cdot));x_\mathrm{lon}^\mathrm{sr},u_\mathrm{brake}(\cdot))) \leq fov, \\
\end{split}
\end{align*}
where $x_\mathrm{lon}^\mathrm{sr} = [0,\speedLimitFOV,\aMaxEgo,\jMaxEgo]^T$ is the worst-case initial state.
%The field of view based velocity limit $\speedLimitFOV$ can change online during driving, e.g., due to changing weather conditions.
Additionally, the ego vehicle has to reduce its velocity to an upcoming speed limit such that it does not brake abruptly. 
Therefore, the speed limit $\speedLimitBraking$ is introduced which ensures that a trajectory for the velocity reduction exists which does not exceed a predefined acceleration threshold $\aAbrupt<0$ and jerk threshold $\jAbrupt<0$. We calculate $\speedLimitBraking$ analogously to $\speedLimitFOV$, where the following two conditions have to be fulfilled:
\begin{align*}
\begin{split}
& x^\mathrm{f}_\mathrm{lon} =  \xi(t^*;x^\mathrm{0}_\mathrm{lon}, u(\cdot)), \\
& \forall t \in [t_0, t^*]: \aEgo(t) \geq \aAbrupt \land \jEgo(t) \geq \jAbrupt, \\
\end{split}
\end{align*}
where $x^\mathrm{0}_\mathrm{lon} = [0, \speedLimitBraking, \aMaxEgo ,\jMaxEgo]^T$ and $x^\mathrm{f}_\mathrm{lon} = [fov, v_\mathrm{sl}^*, 0, 0]^T$ are the worst-case initial state and the final state at which the upcoming speed limit $v_\mathrm{sl}^*$ is reached. 
%The velocity limits $\speedLimitFOV$ and $\speedLimitBraking$ can be calculated offline using binary search. 
Additionally, we assume that the ego vehicle has to consider the maximum allowed velocity $\speedLimitWeather$ based on the road conditions for which we assume that it is externally provided.
%The ego vehicle has also to consider $\speedLimitWeather$ and $\speedLimitType$. 
A vehicle keeps the speed limit given by traffic signs if its velocity is less than $\speedLimitLane$:
\begin{align*}
\laneSpeedLimitPredicate(\kState, v_\mathrm{k}) \iff v_\mathrm{k} \leq \speedLimitLane(\kState).
\end{align*}
The other speed limit predicates are defined analogously.
%\begin{align*}
%\fovSpeedLimitPredicate \iff \vEgo \leq \speedLimitFOV,
%\end{align*}
%\begin{align*}
%\brakingSpeedLimitPredicate \iff \vEgo \leq \speedLimitBraking,
%\end{align*}
%\begin{align*}
%\roadConditionSpeedLimitPredicate \iff \vEgo \leq \speedLimitWeather,
%\end{align*}
%\begin{align*}
%\vehicleTypeSpeedLimitPredicate \iff \vEgo \leq \speedLimitType.
%\end{align*}
A vehicle only needs to consider the minimum required speed limit if its velocity is below $\speedLimitFOV$, $\speedLimitWeather$, and $\speedLimitType$:
\begin{align*}
\begin{split}
&\minSpeedLimitPredicate(\kState, v_\mathrm{k}) \iff \\
&\requiredSpeedLane(\kState) < \min(\speedLimitFOV, \speedLimitWeather, \speedLimitType) \implies \requiredSpeedLane(\kState) \leq v_\mathrm{k}.
\end{split}
\end{align*}
To ensure high traffic throughput on interstates, vehicles must not impede traffic flow:
\begin{align*}
\begin{split}
&\preservesTrafficFlowPredicate(\kState, v_\mathrm{k}) \iff \\ 
& \qquad \neg \slowLeadingVehicle \implies  \mathrm{v\_max}_1(\kState) - v_\mathrm{k}  > \deltaVTrafficFlow,
\end{split}
\end{align*}
where $\mathrm{v\_max}_1(\kState) = \min(\speedLimitBraking, \speedLimitFOV, \speedLimitWeather, v_\mathrm{sl}^{*}(\kState), \speedLimitType(\kState)),$
\begin{align*}
v_\mathrm{sl}^{*}(\kState) =
\begin{cases} 
\speedLimitLane(\kState) & \text{if } \speedLimitLane(\kState) \neq \inf, \\ 
v_\mathrm{su} & \text{otherwise},
\end{cases}	
\end{align*}
$v_\mathrm{su}$ is the suggested velocity which is assumed if no speed limit is present,
\begin{align*}
\begin{split}
&\slowLeadingVehicle(\kState) \iff \\
& \qquad \exists \text{o} \in \{1...N\}: \mathrm{v\_max}_2(\otherState) - \vOther > \deltaVTrafficFlow  \\
& \qquad \land \inSameLaneAsEgoVehiclePredicate(\kState, \otherState) \land \inFrontOfEgoVehiclePredicate(\kState, \otherState),
\end{split}
\end{align*}
$\mathrm{v\_max}_2(\kState) = \min(\speedLimitWeather, v_\mathrm{sl}^{*}(\kState), \speedLimitType(\kState)),$
and $\deltaVTrafficFlow$ is a threshold which indicates if a vehicle drives slowly.
A vehicle is in standstill if its velocity is close to zero:
\begin{align*}
\begin{split}
\inStandStillPredicate(v_\mathrm{k}) \iff -\vError  \leq v_\mathrm{k} \leq \vError, 
\end{split}
\end{align*}
where $v_\mathrm{err}$ captures measurement uncertainties and is close to zero.
This predicate can be used to evaluate if leading vehicles exist which are standing:
\begin{align*}
\begin{split}
& \leadVehicleStandingPredicate(\kState) \iff \\
&  \qquad\exists \text{o} \in \{1...N\}:  \inSameLaneAsEgoVehiclePredicate(\kState, \otherState) \\
& \qquad \land \inFrontOfEgoVehiclePredicate(\kState, \otherState) \land \inStandStillPredicate(\otherState).
\end{split}
\end{align*}
The $\text{k}^\mathrm{th}$ vehicle drives with slightly higher speed than the $\text{p}^\mathrm{th}$ vehicle if the following predicate is true:
\begin{align*}
\begin{split}
& \drivesMaxWithSlightlyHigherSpeedPredicate(\kState, \pState) \iff  \\
& \qquad 0 < \kState - \pState < \velocitySlightlyHigherLimit, 
\end{split}
\end{align*} 
where $\velocitySlightlyHigherLimit$ indicates slightly higher speed. The following proposition indicates that $\text{k}^\mathrm{th}$ vehicle drives faster than any vehicle on its left side:
\begin{align*}
\begin{split}
& \drivesFasterThanVehicleLeftPredicate(v_\mathrm{k}) \iff  \\
& \qquad \exists \text{o} \in \{1...N\}: \vehicleIsLeftPredicate(\otherState) \land v_\mathrm{k} > \vOther.
\end{split}
\end{align*} 
A backward driving vehicle has a negative velocity considering uncertainties:
\begin{align*}
\begin{split}
& \reversesPredicate(v_\mathrm{k}) \iff v_\mathrm{k} < -\vError.
\end{split}
\end{align*}

\subsection{Braking Elements}
The safe distance which has to be kept by the ego vehicle to a leading vehicle is based on the definitions in \cite{Pek2017a} and \cite{Rizaldi2016}:
\begin{align*} 
\safeDistance(\vEgo, \vOther) =  \; \frac{\vOther^2}{2 \aMinOther} - \frac{v_\mathrm{r}^2}{2 \aMinEgo} + \vEgo \timeDelay + \frac{1}{2} \aMaxEgo \timeDelay^2,  
\end{align*} 
where $v_\mathrm{r} = \vEgo + \aMaxEgo \timeDelay$ and \timeDelay denotes the reaction time of the ego vehicle. We assume that other vehicles can brake stronger than the ego vehicle: $\aMinOther < \aMinEgo$.
The safe distance is kept if the relative distance between the vehicles is larger than the safe distance:
\begin{align*}
\begin{split}
&\safeDistanceFrontPredicate(\egoState, \otherState, \vEgo, \vOther) \iff \\ 
& \qquad \safeDistance(\vEgo, \vOther) < (\rearPositionFunction(\otherState) - \frontPositionFunction(\egoState)).
\end{split}
\end{align*}
Unnecessary braking is defined based on the ego vehicle's acceleration and jerk:
\begin{align*}
\begin{split}
&\unnecessaryBraking(\egoState, \aEgo, \jEgo) \iff \\
&\qquad (\neg \velocityReductionNecessary \land \, \aEgo < 0) \implies\\
&\qquad \quad \big(\nexists \text{o} \in \{1 ... N\} : \inFrontOfEgoVehiclePredicate(\kState, \otherState)  \\
&\qquad \quad \land \aEgo < \aAbrupt \land \jEgo < \jAbrupt \big)\\
&\qquad \quad \lor \big(\exists \text{o} \in \{1 ... N\} : \inFrontOfEgoVehiclePredicate(\kState, \otherState) \\
&\qquad \quad \land (\aEgo - \aOther) <  \aAbrupt \land \jEgo  <  \jAbrupt \big),
\end{split}
\end{align*}
where $\velocityReductionNecessary$ indicates if a braking maneuver is necessary due to an externally reason, e.g., another vehicle violates traffic rules so that the ego vehicle has to react to prevent a collision. We assume this information is provided by another system.

\subsection{General Elements}
The $\text{k}^\text{th}$ vehicle is part of a congestion if there is a predefined number of slow driving vehicles in front of it:
\begin{align*}
\begin{split}
& \inCongestion(\kState) \iff  \\
&\qquad \exists^{\numVehiclesCongestion} \, \text{o} \in \{1 ... N\} : \inFrontOfEgoVehiclePredicate(\kState, \otherState) \\
&\qquad \land \inSameLaneAsEgoVehiclePredicate(\otherState) \land \vOther \leq \vCongestion,
\end{split}
\end{align*} 
where \vCongestion indicates slow moving vehicles in a congestion and \numVehiclesCongestion is the minimum required number of vehicles which are part of a congestion within one lane.
A congestion is left of a vehicle if some vehicle on the left side is in a congestion:
\begin{align*}
\begin{split}
& \congestionLeftPredicate(\kState) \iff  \\
&\qquad \exists \text{o} \in \{1 ... N\} : \vehicleIsLeftPredicate(\kState, \otherState)  \\
&\qquad \land \inCongestion(\kState). \\  
\end{split}
\end{align*} 
For a turning vehicle the difference between the vehicle's orientation and the reference path orientation violates a threshold $\uTurnIndicator$:
\begin{align*}
\begin{split}
& \makesUTurnPredicate(\kState) \iff  \\
&\qquad \exists l \in \occupiedLaneletsFunction(\kState): \uTurnIndicator < |\theta_\mathrm{k} - \laneletOrientation(l, s_\mathrm{k})|. \\  
\end{split}
\end{align*} 

